\documentclass[../RFC-HDFG-2026-001.tex]{subfiles}

\begin{document}

%% ======================================================================
%%  Appendix A – VFD/VOL Configuration Profile (Future Work)
%% ======================================================================
\section{VFD/VOL Configuration Profile (Future Work)}
\label{sec:appendix-vfd}

\begin{outofscope}
This entire appendix is out of scope for this RFC.  It is retained as
forward-looking design context, to record that the TOML grammar chosen for
filter configuration (\SectionRef{sec:param-overview}) extends naturally to
VFD and VOL connector configuration without a second syntax.  Nothing
described here---the relaxed VFD/VOL grammar profile, the
\texttt{H5FD\_CONFIG\_*} limits, \texttt{H5FDconfig\_get\_subtable\_str}, or
\texttt{H5FDconfig\_load\_from\_file}---is proposed for implementation by this
RFC, and none of it is present in the HDF5 source tree.  A reader concerned
only with filter configuration may skip the appendix entirely.
\end{outofscope}

\subsection{Additional Grammar Constraints}
\label{sec:vfd-profile-constraints}

VFD and VOL connector configuration strings require richer structure than
filter parameter strings: arbitrary stack depth, arrays (e.g., lists of
endpoint addresses), and file-based delivery that benefits from multi-section
TOML documents.  The following relaxed constraints apply to parameter strings
used in VFD and VOL connector configuration contexts.

\begin{itemize}
\item \textbf{Full TOML v1.0.0 document syntax supported:} Arrays
  (\texttt{[1,2,3]}), multi-line strings, and document-level tables
  (\texttt{[section]}) are all accepted.
\item \textbf{Arbitrary nesting depth:} Inline tables and dotted keys may be
  nested to any depth required to express the VFD stack hierarchy.  There is
  no library-imposed depth limit beyond available stack memory.
\item \textbf{Maximum string length:} \texttt{H5FD\_CONFIG\_STRING\_MAX} (=
  16384 bytes).  A full multi-level VFD stack with a SWMR configuration group
  is unlikely to exceed a few kilobytes; 16\,KB provides ample headroom.
  Strings exceeding this limit are rejected with \texttt{H5E\_BADVALUE}.
\item \textbf{Maximum top-level key count:} \texttt{H5FD\_CONFIG\_MAX\_PARAMS}
  (= 256).  Strings containing more top-level keys are rejected with
  \texttt{H5E\_BADVALUE}.
\item \textbf{Special float values:} \texttt{inf}, \texttt{-inf}, and
  \texttt{nan} are rejected with \texttt{H5E\_BADVALUE}, consistent with the
  filter profile.
\item \textbf{Locale independence:} C-locale numeric conversion applies, as
  in the filter profile.
\end{itemize}

The VFD name is expressed as the sole top-level key, with its configuration
parameters as the value (an inline table or dotted keys).  An underlying VFD
in a stack is expressed as a nested key under the outer VFD's table.  An
empty parameter list is expressed as an empty inline table (\texttt{\{\}}).
For example, a three-level page-buffer~$\to$~encryption~$\to$~sec2 stack:

\begin{minted}{toml}
page_buffer.page_size       = 4096
page_buffer.max_num_pages   = 64
page_buffer.underlying_vfd.encryption_vfd.plaintext_page_size = 4096
page_buffer.underlying_vfd.encryption_vfd.key = "0123456789ABCDEF..."
page_buffer.underlying_vfd.encryption_vfd.underlying_vfd.sec2 = {}
\end{minted}

or equivalently as inline tables:

\begin{minted}{toml}
page_buffer = {
  page_size = 4096, max_num_pages = 64,
  underlying_vfd = { encryption_vfd = {
    plaintext_page_size = 4096, key = "0123456789ABCDEF...",
    underlying_vfd = { sec2 = {} }
  }}
}
\end{minted}

\subsection{Use Case Coverage}
\label{sec:vfd-use-cases}

This subsection demonstrates, use case by use case, that the TOML VFD/VOL
profile is sufficient to express all VFD stack configuration requirements.
The Lifeboat VFD configuration RFC~\cite{lifeboat-vfd-cl} provides a concrete
set of use cases and motivates several of the design choices described hereafter.

\verysubsection{UC-1: Flat scalar VFD parameters}

\textbf{Use case.}  A VFD such as the page buffer VFD requires a small set of
scalar parameters: page size, maximum number of pages, replacement policy.

\textbf{Coverage.}  Standard TOML key\,=\,value pairs cover this
directly.  The grammar provides typed integers, floats, booleans, and quoted
strings without any extension.  Example:

\begin{minted}{toml}
page_buffer = { page_size = 4096, max_num_pages = 64, replacement_policy = 0 }
\end{minted}

No aspect of this use case requires any relaxation beyond ordinary TOML.

\verysubsection{UC-2: Boolean flags}

\textbf{Use case.}  VFD SWMR configuration contains several boolean fields
(e.g., \texttt{presume\_posix\_semantics}, \texttt{flush\_raw\_data}).  The
Lifeboat grammar has no boolean type; it uses integer \texttt{0}/\texttt{1}
as a workaround, requiring comment annotations (\texttt{/* <- Use integer for
bool value */}) to explain the intent.

\textbf{Coverage.}  TOML has a native boolean type:
\texttt{true}/\texttt{false} (lowercase, unquoted).  The intent is
self-documenting and the parser enforces the type without integer-to-boolean
coercion.  Example:

\begin{minted}{toml}
H5F_vfd_swmr_config = {
  presume_posix_semantics = true,
  flush_raw_data          = true,
  generate_updater_files  = false
}
\end{minted}

Native boolean support eliminates the need for integer \texttt{0}/\texttt{1} workarounds.

\verysubsection{UC-3: Binary blob values (e.g., encryption keys)}

\textbf{Use case.}  The encryption VFD requires a raw binary key (32 bytes in
the Lifeboat example).  The Lifeboat grammar introduces a special
\texttt{<binary\_blob>} terminal with a \texttt{--} prefix to distinguish hex
byte sequences from identifiers and integers.

\textbf{Coverage.}  TOML has no native binary type, but a quoted
string carrying a hex-encoded value is fully sufficient.  The VFD callback
receives the string and decodes it.  The convention is explicit and requires no
grammar extension:

\begin{minted}{toml}
encryption_vfd = {
  key = "0123456789ABCDEF0123456789ABCDEF0123456789ABCDEF0123456789ABCDEF"
}
\end{minted}

TOML's quoted string type carries arbitrary content, so no disambiguation
prefix is needed.  The decoding convention (e.g., hex, base64) is specified
by the VFD and validated in its configuration entry point, as with all other
parameter types.

\textbf{Validation responsibility.}  In the Lifeboat grammar, the parser
natively enforces that a \texttt{<binary\_blob>} terminal contains only valid
hex characters.  Under TOML the parser treats the value as an opaque string;
it is therefore the VFD's responsibility to validate the string's content
(character set, byte length, encoding) and return \texttt{H5E\_BADVALUE} if
the value is malformed.  This is the same responsibility pattern already
required of all other parameter types: the TOML grammar validates syntax and
type, while semantic validation (range checks, format checks) belongs to the
VFD.  The shift is in degree, not in kind.

\verysubsection{UC-4: Arbitrary-depth VFD stack nesting}

\textbf{Use case.}  A VFD stack may contain three or more levels (e.g.,
page buffer~$\to$~encryption~$\to$~sec2).  Each VFD must be able to read its
own parameters and forward the remaining configuration to the underlying VFD,
without any VFD needing to know the schema of its neighbors.

\textbf{Coverage.}  TOML inline tables and dotted keys are
recursively defined and support arbitrary nesting depth.  The VFD/VOL profile
removes any library-imposed depth limit.  Each VFD in the stack reads the
keys it recognizes by name and extracts the \texttt{underlying\_vfd} sub-table
to pass to the next level.  Because TOML keys are named rather than
positional, each VFD looks up only what it knows about and ignores everything
else; no VFD needs visibility into any other VFD's schema.

The full three-level example from UC-1:

\begin{minted}{toml}
page_buffer = {
  page_size = 4096, max_num_pages = 64,
  underlying_vfd = { encryption_vfd = {
    plaintext_page_size = 4096, key = "...",
    underlying_vfd = { sec2 = {} }
  }}
}
\end{minted}

The page buffer VFD reads \texttt{page\_size} and \texttt{max\_num\_pages},
then calls \texttt{H5FDconfig\_get\_subtable\_str} to extract the
\texttt{underlying\_vfd} value as a raw TOML string
(\SectionRef{sec:subtable-extraction}), creates a new FAPL, and loads that
string into the FAPL's \texttt{driver\_config\_str} field via
\texttt{H5FDconfig\_load\_str\_into\_fapl}.  The encryption VFD
receives that FAPL, reads its own keys from the extracted string, and
similarly forwards the \texttt{underlying\_vfd.sec2} sub-table to the sec2
configuration.  No VFD ever parses keys belonging to another VFD.

\textbf{API note.}  VFDs do not use the \texttt{H5Z\_class3\_t.set\_config}
callback defined in \SectionRef{sec:class3}; that callback is specific to I/O
filters.  VFD configuration strings are delivered via the
\texttt{driver\_config\_str} field of \texttt{H5FD\_driver\_prop\_t} (defined
in \texttt{H5FDprivate.h}), which is already present in the HDF5 library.
The TOML VFD/VOL profile described here specifies the \emph{format} of that
string; no change to the \texttt{H5FD\_class\_t} struct or the VFD open
callback signature is required by this RFC.

\verysubsection{UC-5: VFD name as the top-level identifier}

\textbf{Use case.}  In the Lifeboat grammar, the VFD name is the outer
identifier of the top-level name-value pair: \texttt{( page\_buffer ( \ldots )
)}.  The configuration string is self-describing: the consumer can determine
which VFD is being configured without out-of-band information.

\textbf{Coverage.}  The same self-describing property is achieved
by convention: the VFD name is the sole top-level key, and its value is the
parameter table.  A configuration string that begins with
\texttt{page\_buffer~=~\{~\ldots~\}} unambiguously identifies the page buffer
VFD as the top-level driver.  The library reads the top-level key name to
determine which VFD to load and configure.

\textbf{Sole top-level key enforcement.}  The library validates that a
VFD/VOL profile configuration string contains \emph{exactly one} top-level
key.  If zero top-level keys are present the string is treated as ``no
configuration'' (equivalent to a \texttt{NULL} string).  If two or more
top-level keys are present (e.g., \texttt{page\_buffer = \{\ldots\}, sec2 =
\{\}}) the library returns \texttt{H5E\_BADVALUE} immediately, before
consulting any VFD.  This rule ensures unambiguous driver identification and
mirrors the Lifeboat grammar's constraint that the top-level is always a
single \texttt{<name\_value\_pair>}.

\verysubsection{UC-6: Empty parameter list}

\textbf{Use case.}  The sec2 VFD requires no configuration parameters.  In
the Lifeboat grammar this is expressed as \texttt{( sec2 () )} --- a VFD name
with an explicit but empty parameter list.

\textbf{Coverage.}  An empty inline table \texttt{sec2 = \{\}} is
the direct and idiomatic equivalent.  The VFD name is present (self-describing),
the parameter list is explicitly empty, and the grammar requires no special
case.  \texttt{NULL} or an empty string is also accepted by the VFD/VOL
profile and treated as ``no parameters'', consistent with the filter profile.

\verysubsection{UC-7: Configuration groups (VFD SWMR)}

\textbf{Use case.}  VFD SWMR requires coordinated configuration across
multiple subsystems: the SWMR engine itself, the page buffer, file space
strategy, and file space page size.  Historically these were configured
independently, risking writer/reader mismatches.  The Lifeboat RFC introduces
a ``configuration group'' concept: a named outer wrapper containing multiple
named sub-configurations that are parsed and applied together.

\textbf{Coverage.}  The configuration group maps directly onto a
TOML table whose value is a list of named sub-tables.  When delivered as an
inline string (e.g., environment variable or API call), nested inline tables
express the group:

\begin{minted}{toml}
vfd_swmr_config_data = {
  H5F_vfd_swmr_config = {
    version = 1, tick_len = 4, max_lag = 7,
    presume_posix_semantics = true, flush_raw_data = true,
    md_file_path = "./md_dir", md_file_name = "md_file"
  },
  page_buffer_config = {
    page_buf_size = 4096, metadata_pages_only = true
  },
  file_space_strategy_config = { persist = false },
  file_space_page_size        = { page_size = 4096 }
}
\end{minted}

When delivered from a file, standard TOML document-level section tables are
the natural form (permitted under the VFD/VOL profile):

\begin{minted}{toml}
[vfd_swmr_config_data.H5F_vfd_swmr_config]
version   = 1
tick_len  = 4
max_lag   = 7
presume_posix_semantics = true
flush_raw_data          = true
md_file_path = "./md_dir"
md_file_name = "md_file"

[vfd_swmr_config_data.page_buffer_config]
page_buf_size       = 4096
metadata_pages_only = true

[vfd_swmr_config_data.file_space_strategy_config]
persist = false

[vfd_swmr_config_data.file_space_page_size]
page_size = 4096
\end{minted}

Both forms parse to the same TOML AST and express identical configuration
data.  However, they are not processed identically at the library level: an
inline string is stored verbatim in the FAPL's \texttt{driver\_config\_str}
and consumed directly by VFDs, while a file containing document-level section
tables is first parsed and normalized to an inline string by the
file-loading helper (\SectionRef{sec:subtable-extraction}) before being
placed in the FAPL.  VFDs always see the inline form; the distinction between
the two authoring styles is transparent to them.

The top-level name (\texttt{vfd\_swmr\_config\_data}) identifies the
configuration group, and each sub-table name identifies the target subsystem.
The writer/reader coordination
problem is solved by the same mechanism: both sides load from the same
configuration string or file, with the library applying writer- and
reader-specific flag overrides (\texttt{writer}, \texttt{create\_file})
programmatically rather than encoding them in the configuration string.

\verysubsection{UC-8: Stack-level delegation without breadth-first parsing}

\textbf{Use case.}  The Lifeboat grammar was designed around a
\emph{breadth-first} parsing model: when a VFD encounters the configuration
string, it parses only its own parameters and receives the sub-expression for
the underlying VFD as a single opaque token, which it passes down the stack
unchanged.  The motivation, as stated in the Lifeboat RFC, is schema
isolation: each VFD in the stack should not need to know anything about the
configuration data for the underlying VFDs.  This is achieved by a
context-sensitive lexer that, when a name-value pair value is expected,
treats a parenthesised sub-expression as a single token rather than
descending into it.  Each VFD parses only the keys at its own level and
forwards the unparsed sub-expression string down the stack.

\textbf{Schema isolation via AST lookup.}  TOML parsers (including the
vendored tomlc17 subset) build a full recursive AST in a single pass, after
which each VFD can locate the keys it recognises by name and extract the
\texttt{underlying\_vfd} sub-table without scanning past unknown structure.
Schema isolation is preserved --- each VFD reads only its own keys --- through
named-key lookup over a complete AST.

\textbf{Implementation.}  Each VFD in the stack receives a TOML configuration
string via \texttt{driver\_config\_str} in its FAPL.  It parses that string,
extracts its own keys by name using \texttt{H5FDconfig\_get\_*} helpers, then
calls \texttt{H5FDconfig\_get\_subtable\_str} (\SectionRef{sec:subtable-extraction})
to extract the raw source bytes of the \texttt{underlying\_vfd} value — no
AST-to-string serialization is required, because the helper works directly on
the original input bytes.  The extracted substring is loaded into a new FAPL's
\texttt{driver\_config\_str} and passed to the next VFD's open callback.
Configuration strings are small (well under the 16\,KB VFD/VOL profile limit),
so the substring extraction overhead is negligible.  Each VFD independently
validates and consumes only its own configuration data, achieving full schema
isolation without any special parsing protocol.

\subsection{Sub-Table String Extraction for VFD Stack Delegation}
\label{sec:subtable-extraction}

VFD stack delegation (\SectionRef{sec:vfd-use-cases}, UC-4 and UC-8) requires
a VFD to pass the configuration of its underlying VFD as a string to that
VFD's configuration entry point.  Because the \texttt{set\_config} entry point
accepts a \texttt{const char~*} string, a VFD cannot pass a parsed
\texttt{toml\_table\_t} pointer directly to the next level.

Standard C TOML parsers, including tomlc17, build an in-memory parse tree but
do not include an emitter to serialize that tree back to a string.  Rather
than adding a full TOML emitter to the vendored library---which would
significantly increase complexity and maintenance burden---the library instead
provides a lightweight helper function that extracts the raw source substring
corresponding to a named sub-table directly from the original input string:

\begin{minted}{c}
herr_t H5FDconfig_get_subtable_str(
    const char  *params,      /* original VFD config string          */
    const char  *key,         /* dotted key path, e.g. "underlying_vfd" */
    char        *buf,         /* output buffer, or NULL for size query   */
    size_t      *buf_size     /* in: capacity; out: required length      */
);
\end{minted}

This function locates the value span of \texttt{key} in \texttt{params} by
scanning the source bytes directly (using the same lexer used by the parser
to identify balanced brace extents), copies the raw substring into
\texttt{buf}, and returns its length via \texttt{buf\_size}.  The extracted
substring is a valid TOML inline table string that the receiving VFD can
parse independently.  This approach:

\begin{itemize}
\item Requires no emitter and no AST traversal beyond balanced-brace matching.
\item Preserves the exact source text, including whitespace and comments,
  avoiding any round-trip fidelity issues.
\item Follows the standard HDF5 size-query pattern: passing \texttt{NULL} for
  \texttt{buf} returns the required buffer size without copying.
\item Returns \texttt{H5E\_NOTFOUND} if the key is absent, and
  \texttt{H5E\_BADVALUE} if the value is not a table (i.e., is a scalar).
\end{itemize}

\textbf{Inline table requirement.}
\texttt{H5FDconfig\_get\_subtable\_str} operates on inline table syntax only.
Balanced-brace scanning works because the value of a named key in an inline
table is always a contiguous, brace-delimited span in the source bytes.
Document-level section tables (\texttt{[section]}) have no such span: the
keys of \texttt{[page\_buffer.underlying\_vfd]} are scattered across the file
with no enclosing braces, so there is no contiguous substring to extract.

This is not a limitation in practice because document-level tables
\emph{never reach \texttt{driver\_config\_str} directly}.  The two delivery
paths are:

\begin{enumerate}
\item \textbf{Inline string} (environment variable, API call): the caller
  supplies an inline table string, which is stored verbatim in
  \texttt{driver\_config\_str}.  \texttt{H5FDconfig\_get\_subtable\_str}
  operates on this string directly.
\item \textbf{File-based delivery}: the application calls a library helper
  such as \texttt{H5FDconfig\_load\_from\_file()} (analogous to the Lifeboat
  \texttt{H5CL\_load\_config\_string\_from\_file()}), which reads the file,
  parses it into a TOML AST (accepting document-level tables in that parse),
  and then emits a single normalized inline table string that it stores in
  the FAPL's \texttt{driver\_config\_str}.  From that point on, all VFDs in
  the stack see only the inline form.
\end{enumerate}

The document-level \texttt{[section]} syntax is therefore a user-facing
authoring convenience.  The library's file-loading helper is the single point
responsible for normalizing it to inline form; \texttt{H5FDconfig\_get\_subtable\_str}
and all downstream VFD configuration logic operate exclusively on inline
strings.  The normalization step requires the file-loading helper to include a
minimal inline-table emitter, scoped only to that one helper and not exposed
to VFD authors.

This function is defined in \texttt{H5FDconfig.c} and declared in
\texttt{H5FDdevelop.h}, placing it in the VFD developer API rather than the
filter developer API (\texttt{H5Zdevelop.h}).  The \texttt{H5FD} prefix
distinguishes it from the filter-scoped \texttt{H5Zconfig\_*} helpers.

\noindent Neither \texttt{H5FDconfig\_get\_subtable\_str} nor the associated
\texttt{H5FDconfig\_load\_from\_file} helper is part of this RFC, and neither is
present in the HDF5 source tree.

\end{document}
